\subsection{An additive-category reduction and its support obstruction}
\label{sec:ringoid}

Proposition~\ref{prop:matrix} puts support and amplitude in one idempotent
equation. We now separate an ordinary additive-category projectivity problem
from a precise obstruction to lifting its support. This first criterion does not eliminate the search for a support section.
Theorem~\ref{thm:canonical-ideal} below replaces that search by a canonical ideal condition, while leaving arbitrary projective $Q$-supports unclassified.

We return to left semimodules. Fix a $Q$-semimodule $L$. Define a preadditive category $\mathcal A_L$
with objects $l\in L$, generated by arrows
\[
 t_{l,k}:l\longrightarrow k\quad(l\leq k),\qquad
 a_{q,r,l}:l\longrightarrow ql\quad(r\in F(q)).
\]
Write $F_{q,q'}$ for a transition of $F$. The relations are
\begin{align*}
 t_{l,l}&=1_l,& t_{k,h}t_{l,k}&=t_{l,h},\\
 a_{q,r+r',l}&=a_{q,r,l}+a_{q,r',l},&
 a_{q',F_{q,q'}(r),l}&=t_{ql,q'l}a_{q,r,l}\quad(q\leq q'),\\
 a_{q,r,k}t_{l,k}&=t_{ql,qk}a_{q,r,l}\quad(l\leq k),&
 a_{q,r,pl}a_{p,s,l}&=a_{qp,\mu_{q,p}(r,s),l},\\
 a_{1_Q,\iota(1),l}&=1_l.&&
\end{align*}
A left $\mathcal A_L$-module is an additive covariant functor
$\mathcal A_L\to\Ab$. This is the usual rings-with-several-objects
formalism \cite{Mitchell,Murfet}; all of the projectivity arguments needed below
are given explicitly. We use projectivity with respect to surjections,
as in \cite[Definition~7.1, p.~18]{BorgerJun}.

\begin{proposition}[fixed support]\label{prop:ringoid}
The category of $T$-semimodules with identified support $L$, and maps
inducing $1_L$, is equivalent to $\mathcal A_L$-modules.
\end{proposition}
\begin{proof}
The transition arrows give the group diagram, and the scalar arrows give
the action homomorphisms in Definition~\ref{def:actiondata}. The transition and naturality relations give the group diagram and its
compatibility with the action; the other relations give scalar
additivity, associativity, and the unit. Scalar addition across
different supports follows by transporting both summands to their join
and then using additivity there. Module addition across different fibers
follows in the same way from transition naturality. Thus
Theorem~\ref{thm:general} reconstructs the semimodule, and its morphism
statement gives the asserted equivalence. The zero object of this
additive category has the zero group at every $l$; its underlying
$T$-semimodule is $L$, which need not be the zero semimodule.
\end{proof}

\begin{proposition}[finite generation]\label{prop:ringoid-finite}
A $T$-semimodule $M$, with support $L$ and amplitude functor $G$, is
finitely generated if and only if $L$ is finitely generated over $Q$
and $G$ is finitely generated over $\mathcal A_L$.
\end{proposition}
\begin{proof}
The supports of finite semimodule generators generate $L$. The same
elements generate $G$: each scalar multiple is a scalar arrow, and a
sum is formed by transitions to a common fiber followed by group addition.

Conversely choose finite $\mathcal A_L$-generators of $G$ and finite
$Q$-generators $l_i$ of $L$. Use the former elements and the fiber zeros
$0_{l_i}$ as semimodule generators. All fiber zeros are generated, because
zero amplitudes in $F(q)$ move $0_{l_i}$ to $0_{ql_i}$ and addition joins
their supports. A transition to a larger fiber is realized by adding its
fiber zero. The scalar $-1_T$ gives the group inverse at the same
support. Therefore every additive-category expression in the selected
generators is an expression in the generated semimodule.
\end{proof}

Choose any finite free cover
\[
 f:T^n\twoheadrightarrow M,\qquad
 (s_j)\longmapsto\sum_j s_jm_j,
 \qquad m_j\in G(l_j).
\]
Its support map is
\[
 \rho:Q^n\twoheadrightarrow L,\qquad
 (q_j)\longmapsto\bigvee_jq_jl_j.
\]
For a $Q$-linear section $\sigma:L\to Q^n$, let
\[
 N_\sigma(l)=\bigoplus_{j=1}^nF\bigl(\sigma(l)_j\bigr).
\]
The free-module transitions and actions restrict along $\sigma$ to make
$N_\sigma$ an $\mathcal A_L$-module. Restriction of $f$ is an
$\mathcal A_L$-linear map $c_\sigma:N_\sigma\to G$. Define
\[
 C_\sigma=\operatorname{coker}(c_\sigma)
 \quad\text{in }\mathcal A_L\text{-}\Mod.
\]

\begin{theorem}[support-section criterion]\label{thm:ringoid-projective}
A finitely generated $T$-semimodule $M$ is projective if and only if
\begin{enumerate}
 \item[\textup{(i)}] $L$ is a finite projective $Q$-semimodule;
 \item[\textup{(ii)}] $G$ is a finitely generated projective $\mathcal A_L$-module;
 \item[\textup{(iii)}] some $Q$-linear section $\sigma$ of $\rho$ satisfies $C_\sigma=0$.
\end{enumerate}
The criterion is independent of the selected finite free cover. For a
fixed $\sigma$, condition \textup{(iii)} is equivalent to the finitely
many linear image conditions
\[
 m_j\in\im\bigl(c_\sigma(l_j)\bigr)\qquad(1\leq j\leq n).
\]
The existence of $\sigma$ already implies \textup{(i)}; listing it
separately isolates the first support-theoretic test.
\end{theorem}
\begin{proof}
Suppose $M$ is projective. Any finite free cover $f$ splits. The support
$\sigma$ of a splitting satisfies $\rho\sigma=1_L$, so $L$ is a
finite free retract. On group fibers the splitting lands in
$N_\sigma(l)$ and splits $c_\sigma$. Hence $C_\sigma=0$.

Let $E\twoheadrightarrow G$ be an epimorphism of $\mathcal A_L$-modules.
It is pointwise surjective, so reconstruction gives a surjective
$T$-map onto $M$ with support $1_L$. Projectivity of $M$ splits it.
The support of this splitting must also be $1_L$, and therefore the
splitting is $\mathcal A_L$-linear. Every epimorphism onto $G$ splits,
which proves projectivity of $G$. Its finite generation follows from
Proposition~\ref{prop:ringoid-finite}.

Conversely suppose \textup{(ii)} and \textup{(iii)} hold.
Projectivity of $G$ splits the pointwise epimorphism $c_\sigma$ by an
$\mathcal A_L$-map $u:G\to N_\sigma$. Reconstruct this as a $T$-map
with support $1_L$, and then include each restricted free fiber into
$T^n$ with support map $\sigma$. The resulting $h:M\to T^n$
satisfies $fh=1_M$ on groups and supports. Thus $M$ is a finite free
retract. Finally, $\im c_\sigma$ is an $\mathcal A_L$-submodule and
the $m_j$ generate $G$, proving the finite image test. Projectivity of
$G$, not arbitrary choices of pointwise lifts, supplies the coherence
of the section.
\end{proof}

\begin{proposition}[a nonzero support obstruction]\label{prop:ringoid-obstruction}
For $T=R_0$ with $R\ne0$, projectivity of the support and of its
amplitude module does not imply projectivity over $T$.
\end{proposition}
\begin{proof}
Take $L=\B^2$ with atoms $a,b$ and top $t=a\vee b$. Here
$\mathcal A_L$ is the $R$-linear incidence category. The diagram
$G(t)=R$, $G(l)=0$ for $l\ne t$, is its projective representable
born at $t$. Its semimodule is generated by $0_a,0_b$ and $1\in G(t)$.

For any finite free cover and any support section,
$\sigma(t)=\sigma(a)\vee\sigma(b)$. Every active free coordinate
at $\sigma(t)$ is therefore active at $\sigma(a)$ or $\sigma(b)$.
The free group at $\sigma(t)$ is the sum of the transition images
of these two lower groups. Their images under $f$ are zero, since
$G(a)=G(b)=0$. Thus $c_\sigma(t)=0$ and
$C_\sigma(t)=R\ne0$ for every $\sigma$. Theorem~\ref{thm:ringoid-projective}
excludes projectivity.
\end{proof}

For $G_{\operatorname{Id}(C)}(R)$, the action category is again the
$R$-linear incidence category of its support: all proper scalar group
fibers vanish. Its finite projective diagrams may have birth labels at
arbitrary objects. Theorem~\ref{thm:finite-lattice-projectives}
identifies exactly which births pass the additional support test,
namely the admissible ideals. The next subsection replaces the presentation-dependent support-section obstruction by a canonical ideal for arbitrary $(Q,F)$.
